\documentclass[12pt,a4paper]{article}

% ===== PACKAGES =====
\usepackage[utf8]{inputenc}
\usepackage[vietnamese]{babel}
\usepackage{amsmath, amssymb, amsfonts}
\usepackage{geometry}
\usepackage{enumitem}
\usepackage{array}
\usepackage{booktabs}
\usepackage{fancyhdr}
\usepackage{titlesec}
\usepackage{xcolor}
\usepackage{tcolorbox}

\geometry{a4paper, margin=2.5cm}

% ===== CUSTOM STYLES =====
\definecolor{mainblue}{RGB}{0,90,160}
\definecolor{darkgreen}{RGB}{0,120,60}
\definecolor{resultbg}{RGB}{230,245,255}

\newtcolorbox{resultbox}{
  colback=resultbg, colframe=mainblue,
  boxrule=0.8pt, arc=3pt,
  left=8pt, right=8pt, top=6pt, bottom=6pt,
  fonttitle=\bfseries, title=Kết quả
}

\newtcolorbox{notebox}[1][]{
  colback=yellow!8, colframe=orange!70!black,
  boxrule=0.6pt, arc=3pt,
  left=8pt, right=8pt, top=6pt, bottom=6pt,
  fonttitle=\bfseries, title={Lưu ý}, #1
}

\pagestyle{fancy}
\fancyhf{}
\fancyhead[L]{\textit{Phân tích \& Thiết kế Thuật toán}}
\fancyhead[R]{\textit{GV: Huỳnh Thị Thanh Thương}}
\fancyfoot[C]{\thepage}

\title{\textbf{Bài giải Bài 1: Sắp xếp các hàm theo bậc tăng trưởng}\\[6pt]
\large (Order of Growth -- Có lập luận/giải thích)}
\author{}
\date{}

\begin{document}
\maketitle
\thispagestyle{fancy}

% ===== PHƯƠNG PHÁP =====
\section*{Phương pháp giải (theo hướng dẫn của giảng viên)}

Theo cô Thương hướng dẫn trong buổi 5, khi so sánh bậc tăng trưởng của các hàm:

\begin{enumerate}[label=\textbf{(\arabic*)}]
  \item \textbf{Không đoán mò} --- phải có biến đổi toán học làm cơ sở.
  \item \textbf{Đưa về cùng cơ số} --- nếu hai hàm đều ở dạng mũ nhưng khác cơ số, chuyển hết về cơ số nhỏ nhất (thường là $2$) rồi so phần mũ.
  \item \textbf{Phân lớp (phân nhóm)} --- quy mỗi hàm về lớp độ phức tạp quen thuộc ($\Theta$ tương ứng), rồi chỉ cần so sánh giữa các lớp.
  \item \textbf{Hệ số bỏ đi} --- trong cùng lớp, hệ số hằng là không đáng kể; chỉ giữ lại phần ``bậc'' cao nhất.
  \item \textbf{Kết luận}: ghi rõ hàm thuộc lớp $O/\Theta$ nào (được \textbf{0.25 điểm/hàm} khi thi), không cần ghi đúng công thức gốc.
\end{enumerate}

\begin{notebox}
Thứ tự các lớp cơ bản cần nhớ:
\[
  \underbrace{\Theta(1)}_{\text{hằng}} \prec
  \underbrace{\Theta(\log n)}_{\text{logarit}} \prec
  \underbrace{\Theta(n^c)}_{\text{đa thức}} \prec
  \underbrace{\Theta(c^n)}_{\text{hàm mũ}} \prec
  \underbrace{\Theta(n!)}_{\text{giai thừa}} \prec
  \underbrace{\Theta(c^{c^n})}_{\text{siêu mũ}}
\]
Trong cùng lớp đa thức: so bậc. Trong cùng lớp mũ: đưa về cùng cơ số, so phần mũ.
\end{notebox}

% =========================================
% NHÓM 1
% =========================================
\newpage
\section{Nhóm 1 (Group 1 --- 7 hàm)}

\textbf{Đề bài:} Sắp xếp các hàm sau theo bậc tăng trưởng từ thấp đến cao:
\[
  (n-2)!, \quad 5\lg(n+100)^{10}, \quad 2^{2n}, \quad 0.001n^4 + 3n^3 + 1, \quad \ln^2 n, \quad \sqrt[3]{n}, \quad 3^n.
\]

\subsection{Bước 1: Rút gọn và phân lớp từng hàm}

\begin{enumerate}[label=\textbf{\arabic*.}]

  \item $5\lg(n+100)^{10}$

  Theo tính chất logarit: $\lg(n+100)^{10} = 10\lg(n+100)$.

  Vậy: $5\lg(n+100)^{10} = 50\lg(n+100)$.

  Khi $n$ đủ lớn: $\lg(n+100) \approx \lg n$ (sai biệt không đáng kể).

  $\Longrightarrow$ Hàm này thuộc lớp $\Theta(\log n)$.

  \item $\ln^2 n$

  Tức là $(\ln n)^2$. Đây là hàm logarit bình phương.

  $\Longrightarrow$ Thuộc lớp $\Theta(\log^2 n)$.

  \item $\sqrt[3]{n}$

  Viết lại dưới dạng số mũ: $\sqrt[3]{n} = n^{1/3}$.

  $\Longrightarrow$ Thuộc lớp $\Theta(n^{1/3})$ --- đây là lớp \textbf{đa thức}, bậc $1/3$.

  \item $0.001n^4 + 3n^3 + 1$

  Là đa thức bậc 4 theo $n$. Bỏ hệ số, bỏ các bậc thấp hơn, chỉ giữ bậc cao nhất.

  $\Longrightarrow$ Thuộc lớp $\Theta(n^4)$.

  \item $3^n$

  \textbf{Chuyển về cơ số 2} (theo phương pháp cô hướng dẫn):
  \[
    3^n = \bigl(2^{\lg 3}\bigr)^n = 2^{n\lg 3} \approx 2^{1{,}585\,n}.
  \]
  Phần mũ xấp xỉ $O(n)$. $\Longrightarrow$ Thuộc lớp hàm mũ $\Theta(3^n)$.

  \item $2^{2n}$

  Rút gọn:
  \[
    2^{2n} = \bigl(2^2\bigr)^n = 4^n = 2^{2n}.
  \]
  \textbf{Chuyển về cơ số 2}: phần mũ là $2n$.

  So với $3^n$ (phần mũ $\approx 1{,}585\,n$): ta có $2n > 1{,}585\,n$ nên $4^n > 3^n$.

  $\Longrightarrow$ Thuộc lớp hàm mũ $\Theta(4^n)$.

  \item $(n-2)!$

  Đây là hàm giai thừa. Khi $n$ đủ lớn, $(n-2)!$ tăng nhanh hơn mọi hàm mũ.

  $\Longrightarrow$ Thuộc lớp $\Theta((n-2)!) = \Theta(n!)$.

\end{enumerate}

\subsection{Bước 2: Lập luận so sánh giữa các lớp}

\begin{enumerate}[label=\textbf{(\alph*)}]
  \item \textbf{Logarit $\prec$ Logarit bình phương:}
    $\log n \prec \log^2 n$ (bậc 1 < bậc 2 của log).

    $\Rightarrow 5\lg(n+100)^{10} \prec \ln^2 n$.

  \item \textbf{Logarit $\prec$ Đa thức:}
    Mọi hàm logarit (dù lũy thừa bao nhiêu) luôn tăng chậm hơn bất kỳ hàm đa thức nào (dù bậc nhỏ tới đâu).

    $\Rightarrow \ln^2 n \prec \sqrt[3]{n}$.

  \item \textbf{Đa thức bậc thấp $\prec$ Đa thức bậc cao:}
    So bậc: $1/3 < 4$.

    $\Rightarrow \sqrt[3]{n} \prec 0.001n^4 + 3n^3 + 1$.

  \item \textbf{Đa thức $\prec$ Hàm mũ:}
    Hàm đa thức (dù bậc cao tới đâu) luôn tăng chậm hơn hàm mũ.

    $\Rightarrow n^4 \prec 3^n$.

  \item \textbf{Hàm mũ cơ số nhỏ $\prec$ Hàm mũ cơ số lớn:}
    Đưa về cùng cơ số 2: phần mũ $1{,}585\,n < 2n$ nên $3^n \prec 4^n$.

    $\Rightarrow 3^n \prec 2^{2n}$.

  \item \textbf{Hàm mũ $\prec$ Giai thừa:}
    Hàm giai thừa tăng nhanh hơn bất kỳ hàm mũ nào.

    $\Rightarrow 2^{2n} \prec (n-2)!$.
\end{enumerate}

\begin{resultbox}
\textbf{Kết quả sắp xếp Nhóm 1} (từ thấp đến cao):
\[
  \boxed{
    5\lg(n+100)^{10}
    \;\prec\; \ln^2 n
    \;\prec\; \sqrt[3]{n}
    \;\prec\; 0.001n^4 + 3n^3 + 1
    \;\prec\; 3^n
    \;\prec\; 2^{2n}
    \;\prec\; (n-2)!
  }
\]
\end{resultbox}

% =========================================
% NHÓM 2
% =========================================
\newpage
\section{Nhóm 2 (Group 2 --- 30 hàm, bài 3-3 CLRS)}

\textbf{Đề bài:} Sắp xếp 30 hàm theo bậc tăng trưởng và phân thành các lớp tương đương
($f \equiv g$ khi $f = \Theta(g)$).

\textbf{Quy ước:} $\lg n = \log_2 n$, \; $\ln n = \log_e n$, \; $\lg^k n = (\lg n)^k$, \; $\lg\lg n = \lg(\lg n)$.

\bigskip

Chia 30 hàm thành \textbf{10 lớp tương đương}, sắp từ nhỏ đến lớn.

% ----- LỚP 1 -----
\subsection{Lớp 1: Hằng số --- $\Theta(1)$}

\textbf{Các hàm:} $1$ \;\; và \;\; $n^{1/\lg n}$.

\textbf{Lập luận cho $n^{1/\lg n}$:}

Đặt $y = n^{1/\lg n}$. Lấy $\lg$ hai vế:
\[
  \lg y = \frac{1}{\lg n} \cdot \lg n = 1
  \quad\Longrightarrow\quad y = 2^1 = 2.
\]
Vậy $n^{1/\lg n} = 2$ (hằng số) $\in \Theta(1)$.

\begin{resultbox}
$1 \;\equiv\; n^{1/\lg n}$
\end{resultbox}

% ----- LỚP 2 -----
\subsection{Lớp 2: Logarit lặp --- $\Theta(\lg^* n)$ và lân cận}

\textbf{Các hàm:} $\lg(\lg^* n)$, \;\; $\lg^* n$, \;\; $\lg^*(\lg n)$, \;\; $2^{\lg^* n}$.

\textbf{Lập luận:}
\begin{itemize}
  \item $\lg^*(\lg n) = \lg^* n - 1 \in \Theta(\lg^* n)$.
    Nên $\lg^* n \equiv \lg^*(\lg n)$.
  \item $\lg(\lg^* n)$: lấy log của hàm log lặp $\Rightarrow$ tăng chậm hơn bản thân $\lg^* n$.
  \item $2^{\lg^* n}$: hàm mũ của $\lg^* n$ $\Rightarrow$ tăng nhanh hơn $\lg^* n$.
\end{itemize}

\begin{resultbox}
$\lg(\lg^* n)
\;\prec\; \lg^* n \equiv \lg^*(\lg n)
\;\prec\; 2^{\lg^* n}$
\end{resultbox}

% ----- LỚP 3 -----
\subsection{Lớp 3: Log của log --- $\Theta(\lg\lg n)$}

\textbf{Hàm:} $\ln\ln n$.

\textbf{Lập luận:} $\ln\ln n = \Theta(\lg\lg n)$ (chỉ khác hệ số do đổi cơ số log).
Lớn hơn $\lg^* n$ nhưng nhỏ hơn các hàm logarit chuẩn.

% ----- LỚP 4 -----
\subsection{Lớp 4: Logarit --- $\Theta(\log^k n)$}

\textbf{Các hàm:} $\sqrt{\lg n}$, \;\; $\ln n$, \;\; $\lg^2 n$.

\textbf{Lập luận:}
\begin{itemize}
  \item $\sqrt{\lg n} = (\lg n)^{1/2}$ --- logarit bậc $1/2$.
  \item $\ln n = \Theta(\lg n)$ --- logarit bậc $1$ (chỉ khác hệ số cơ số).
  \item $\lg^2 n = (\lg n)^2$ --- logarit bậc $2$.
\end{itemize}
So bậc: $1/2 < 1 < 2$.

\begin{resultbox}
$\sqrt{\lg n} \;\prec\; \ln n \;\prec\; \lg^2 n$
\end{resultbox}

% ----- LỚP 5 -----
\subsection{Lớp 5: Siêu logarit / Tiền đa thức}

\textbf{Hàm:} $2^{\sqrt{2\lg n}}$.

\textbf{Lập luận:} Lấy $\lg$ hai vế: $\lg\bigl(2^{\sqrt{2\lg n}}\bigr) = \sqrt{2\lg n}$.

So với $\lg^2 n = (\lg n)^2$: khi $n$ lớn, $\sqrt{2\lg n} \gg (\lg n)^2$? \textbf{Không} ---
thực tế $\sqrt{\lg n}$ tăng chậm hơn $(\lg n)^2$. Tuy nhiên, hàm $2^{\sqrt{2\lg n}}$ tăng
\textbf{nhanh hơn} mọi $\log^k n$ vì phần mũ đang tăng (dù chậm).

So với đa thức $\sqrt{n}$: $\lg(\sqrt{n}) = \frac{1}{2}\lg n \gg \sqrt{2\lg n}$ khi $n$ lớn.

$\Longrightarrow$ Nằm giữa lớp logarit và lớp đa thức.

% ----- LỚP 6 -----
\subsection{Lớp 6: Đa thức --- $\Theta(n^c)$}

\textbf{Các hàm:} $(\sqrt{2})^{\lg n}$, \; $n$, \; $2^{\lg n}$, \; $\lg(n!)$, \; $n\lg n$, \; $n^2$, \; $4^{\lg n}$, \; $n^3$.

\textbf{Lập luận chi tiết:}

\begin{enumerate}[label=(\roman*)]
  \item $(\sqrt{2})^{\lg n}$:

  Chuyển về cơ số 2:
  $(\sqrt{2})^{\lg n} = (2^{1/2})^{\lg n} = 2^{\frac{1}{2}\lg n} = (2^{\lg n})^{1/2} = n^{1/2} = \sqrt{n}$.

  $\Longrightarrow \Theta(n^{1/2})$.

  \item $2^{\lg n}$:

  Theo định nghĩa: $2^{\lg n} = 2^{\log_2 n} = n$.

  $\Longrightarrow \Theta(n)$. Vậy $n \equiv 2^{\lg n}$.

  \item $\lg(n!)$:

  Theo công thức Stirling (như cô nhắc trong buổi 5):
  \[
    \lg(n!) = \lg 1 + \lg 2 + \cdots + \lg n \leq n\lg n
    \quad\Longrightarrow\quad \lg(n!) \in O(n\lg n).
  \]
  Chiều ngược lại cũng chứng minh được $\lg(n!) \in \Omega(n\lg n)$.

  $\Longrightarrow \lg(n!) \in \Theta(n\lg n)$. Vậy $n\lg n \equiv \lg(n!)$.

  \item $4^{\lg n}$:

  Chuyển về cơ số 2:
  $4^{\lg n} = (2^2)^{\lg n} = 2^{2\lg n} = (2^{\lg n})^2 = n^2$.

  $\Longrightarrow \Theta(n^2)$. Vậy $n^2 \equiv 4^{\lg n}$.
\end{enumerate}

So bậc các đa thức: $1/2 < 1 < n\lg n < 2 < 3$.

\begin{resultbox}
$(\sqrt{2})^{\lg n}
\;\prec\; n \equiv 2^{\lg n}
\;\prec\; n\lg n \equiv \lg(n!)
\;\prec\; n^2 \equiv 4^{\lg n}
\;\prec\; n^3$
\end{resultbox}

% ----- LỚP 7 -----
\subsection{Lớp 7: Siêu đa thức / Tiền hàm mũ}

\textbf{Các hàm:} $(\lg n)!$, \;\; $n^{\lg\lg n}$, \;\; $(\lg n)^{\lg n}$.

\textbf{Lập luận:}

\begin{enumerate}[label=(\roman*)]
  \item $(\lg n)^{\lg n}$:

  Lấy $\lg$ hai vế: $\lg\bigl((\lg n)^{\lg n}\bigr) = \lg n \cdot \lg\lg n$.

  \item $n^{\lg\lg n}$:

  Lấy $\lg$ hai vế: $\lg\bigl(n^{\lg\lg n}\bigr) = \lg\lg n \cdot \lg n$.

  $\Longrightarrow$ Hai hàm có cùng log $\Rightarrow$ $n^{\lg\lg n} \equiv (\lg n)^{\lg n}$.

  \item $(\lg n)!$:

  Áp dụng Stirling: $\lg\bigl((\lg n)!\bigr) \approx \lg n \cdot \lg\lg n$, nhưng có hệ số nhỏ hơn
  khi phân tích kỹ. Nên $(\lg n)! \prec n^{\lg\lg n}$.
\end{enumerate}

\begin{resultbox}
$(\lg n)! \;\prec\; n^{\lg\lg n} \equiv (\lg n)^{\lg n}$
\end{resultbox}

% ----- LỚP 8 -----
\subsection{Lớp 8: Hàm mũ --- $\Theta(c^n)$}

\textbf{Các hàm:} $(3/2)^n$, \;\; $2^n$, \;\; $n \cdot 2^n$, \;\; $e^n$.

\textbf{Lập luận --- đưa về cùng cơ số 2:}

\begin{enumerate}[label=(\roman*)]
  \item $(3/2)^n = 2^{n\lg(3/2)} \approx 2^{0{,}585\,n}$. Phần mũ: $O(n)$, hệ số $\approx 0{,}585$.

  \item $2^n$: Phần mũ đúng $n$, hệ số $= 1$.

  \item $n \cdot 2^n$: Lấy $\lg$: $n + \lg n$. Phần chính vẫn là $n$ nhưng có thêm $\lg n$ $\Rightarrow$ lớn hơn $2^n$ một chút.

  \item $e^n = 2^{n\lg e} \approx 2^{1{,}443\,n}$. Phần mũ: hệ số $\approx 1{,}443 > 1$.
\end{enumerate}

So phần mũ (hệ số): $0{,}585 < 1 < 1 + \frac{\lg n}{n} < 1{,}443$.

\begin{resultbox}
$(3/2)^n \;\prec\; 2^n \;\prec\; n \cdot 2^n \;\prec\; e^n$
\end{resultbox}

% ----- LỚP 9 -----
\subsection{Lớp 9: Giai thừa --- $\Theta(n!)$}

\textbf{Các hàm:} $n!$, \;\; $(n+1)!$.

\textbf{Lập luận:}
\begin{itemize}
  \item $n!$ tăng nhanh hơn $e^n$ vì $\lg(n!) = \Theta(n\lg n)$ mà $\lg(e^n) = n\lg e = \Theta(n)$, và $n\lg n \gg n$.
  \item $(n+1)! = (n+1) \cdot n!$ --- rõ ràng tăng nhanh hơn $n!$.
\end{itemize}

\begin{resultbox}
$n! \;\prec\; (n+1)!$
\end{resultbox}

% ----- LỚP 10 -----
\subsection{Lớp 10: Siêu hàm mũ (Double Exponential)}

\textbf{Các hàm:} $2^{2^n}$, \;\; $2^{2^{n+1}}$.

\textbf{Lập luận:}
\[
  2^{2^{n+1}} = 2^{2 \cdot 2^n} = \bigl(2^{2^n}\bigr)^2.
\]
Bình phương của $2^{2^n}$ rõ ràng lớn hơn rất nhiều so với bản thân nó.

\begin{resultbox}
$2^{2^n} \;\prec\; 2^{2^{n+1}}$
\end{resultbox}

% ===== TỔNG KẾT =====
\newpage
\section*{Tổng kết: Thứ tự hoàn chỉnh 30 hàm (Nhóm 2)}

\begin{resultbox}
\small
\[
\begin{aligned}
  & \underbrace{1 \equiv n^{1/\lg n}}_{\Theta(1)}
    \;\prec\; \lg(\lg^* n)
    \;\prec\; \underbrace{\lg^* n \equiv \lg^*(\lg n)}_{\Theta(\lg^* n)}
    \;\prec\; 2^{\lg^* n} \\[6pt]
  & \;\prec\; \underbrace{\ln\ln n}_{\Theta(\lg\lg n)}
    \;\prec\; \underbrace{\sqrt{\lg n}}_{\Theta(\lg^{1/2} n)}
    \;\prec\; \underbrace{\ln n}_{\Theta(\lg n)}
    \;\prec\; \underbrace{\lg^2 n}_{\Theta(\lg^2 n)}
    \;\prec\; 2^{\sqrt{2\lg n}} \\[6pt]
  & \;\prec\; \underbrace{(\sqrt{2})^{\lg n}}_{= \sqrt{n}}
    \;\prec\; \underbrace{n \equiv 2^{\lg n}}_{\Theta(n)}
    \;\prec\; \underbrace{n\lg n \equiv \lg(n!)}_{\Theta(n\lg n)}
    \;\prec\; \underbrace{n^2 \equiv 4^{\lg n}}_{\Theta(n^2)}
    \;\prec\; n^3 \\[6pt]
  & \;\prec\; \underbrace{(\lg n)!}_{} 
    \;\prec\; \underbrace{n^{\lg\lg n} \equiv (\lg n)^{\lg n}}_{} \\[6pt]
  & \;\prec\; \underbrace{(3/2)^n}_{}
    \;\prec\; 2^n
    \;\prec\; n \cdot 2^n
    \;\prec\; e^n \\[6pt]
  & \;\prec\; n!
    \;\prec\; (n+1)!
    \;\prec\; 2^{2^n}
    \;\prec\; 2^{2^{n+1}}
\end{aligned}
\]
\end{resultbox}

\bigskip

\begin{notebox}[title={Mẹo khi đi thi (theo cô Thương)}]
\begin{itemize}
  \item Chỉ cần ghi hàm thuộc lớp $\Theta$ nào là đã được \textbf{0.25--0.5 điểm/hàm}.
  \item Không cần ghi đúng công thức gốc --- ví dụ chỉ cần viết ``$\sum_{i=1}^{n} i^2 \in \Theta(n^3)$''.
  \item Khi so sánh hai hàm mũ: \textbf{chuyển về cùng cơ số 2}, rồi so phần mũ.
  \item Nếu không chắc đúng/sai $\Rightarrow$ dùng \textbf{phản chứng}: giả sử đúng, biến đổi tới mâu thuẫn.
  \item Khi dùng giới hạn: $\displaystyle\lim_{n\to\infty}\frac{f(n)}{g(n)} = 0 \Rightarrow f \in O(g)$; $= c \neq 0 \Rightarrow f = \Theta(g)$; $= \infty \Rightarrow g \in O(f)$.
\end{itemize}
\end{notebox}


